%% secoes/resultados/04-apgar.tex — incluído por main.tex

%%---------------------------------------------------------------------%%
\subsection{Apgar baixo no quinto minuto}\label{subsec:res-apgar}
%%---------------------------------------------------------------------%%

A Tabela~\ref{tab:apgar-dist} mostra a distribuição do Apgar baixo no quinto minuto nas oito populações. No recorte geral, a prevalência basal do desfecho foi de 2,19\%.
Seis dos sete segmentos vulneráveis ficaram acima desse valor: população quilombola (3,26\%), mães menores de 15 anos (3,13\%), mães com baixa escolaridade (2,64\%), população rural (2,56\%), população indígena (2,30\%) e municípios de até 10 mil habitantes (2,29\%). Apenas as beneficiárias do Bolsa Família ficaram ligeiramente abaixo (2,17\%).

\begin{table}[H]
\caption{Distribuição do Apgar baixo no quinto minuto por população analisada.}\label{tab:apgar-dist}%
\begin{tabularx}{\textwidth}{@{}Yrrrr@{}}
\toprule
População & Total & \makecell{Total\\(\%)} & \makecell{Nº de\\ocorrências} & \makecell{Ocorrências\\(\%)}\\
\midrule
Recorte geral & 5.040.358 & 100,00 & 110.267 & 2,19\\
Municípios até 10 mil hab. & 404.920 & 8,03 & 9.259 & 2,29\\
Rural & 1.046.564 & 20,76 & 26.812 & 2,56\\
Indígena & 37.446 & 0,74 & 863 & 2,30\\
Quilombola & 21.628 & 0,43 & 706 & 3,26\\
Mães com baixa escolaridade & 194.344 & 3,86 & 5.140 & 2,64\\
Mães menores de 15 anos & 25.627 & 0,51 & 802 & 3,13\\
Beneficiárias do Bolsa Família & 1.167.427 & 23,16 & 25.324 & 2,17\\
\bottomrule
\end{tabularx}
%\fontetabela
\end{table}


A Tabela~\ref{tab:apgar} apresenta os cinco subgrupos de maior WRAcc para o Apgar baixo no quinto minuto nos oito recortes analisados.
No Apgar baixo no quinto minuto, os itens mais recorrentes descrevem o recém-nascido e o parto. O sexo masculino aparece em seis dos oito recortes, ausente apenas nas populações indígena e quilombola; isolado, está entre 8\% e 10\% acima da basal (mães com baixa escolaridade e recorte geral). O parto normal aparece em quatro recortes (recorte geral, municípios de até 10 mil habitantes, rural e mães menores de 15 anos), e a conjunção de parto normal e primeira gestação se repete em três deles com risco semelhante: 36\% no recorte geral, 40\% nos municípios de até 10 mil habitantes e 36\% no rural. A primeira gestação isolada está 30\% acima da basal na população indígena e 25\% entre as beneficiárias do Bolsa Família.


\begin{table}[H]
\centering
\caption{Subgrupos para Apgar baixo no quinto minuto (prevalência basal: 2,19\%).}
\label{tab:apgar}

\scriptsize
\setlength{\tabcolsep}{4pt}
\renewcommand{\arraystretch}{1.04}

\begin{tabularx}{\textwidth}{@{}c Y r r r r r@{}}
\toprule
\textbf{\#} & \textbf{Subgrupo} & \textbf{Total} & \textbf{\makecell[r]{Total\\(\%)}} & \textbf{\makecell[r]{Nº de\\ocorr.}} & \textbf{\makecell[r]{Ocorr.\\(\%)}} & \textbf{Lift} \\
\midrule
\secaotabela{Recorte geral}
1 & sexo: M & 2.574.637 & 51,1 & 61.992 & 56,2 & 1,10 \\
2 & parto normal; primeira gestação: sim & 659.766 & 13,1 & 19.615 & 17,8 & 1,36 \\
3 & sexo: M; parto normal; raça/cor da mãe: preta/parda & 915.585 & 18,2 & 24.513 & 22,2 & 1,22 \\
4 & calçamento do domicílio: não existe; sexo: M & 978.981 & 19,4 & 25.790 & 23,4 & 1,20 \\
5 & sexo: M; primeira gestação: sim & 670.630 & 13,3 & 18.989 & 17,2 & 1,29 \\

\secaotabela{Municípios de até 10 mil habitantes}
6 & sexo: M; parto normal & 99.644 & 24,6 & 2.817 & 30,4 & 1,24 \\
7 & consultas: 4--6 & 114.133 & 28,2 & 3.102 & 33,5 & 1,19 \\
8 & parto normal; primeira gestação: sim & 51.947 & 12,8 & 1.668 & 18,0 & 1,40 \\
9 & material do piso: cerâmica/pedra/carpete; sexo: M & 120.069 & 29,7 & 3.202 & 34,6 & 1,17 \\
10 & parto normal & 197.586 & 48,8 & 4.902 & 52,9 & 1,08 \\

\secaotabela{População rural}
11 & sexo: M; parto normal & 314.978 & 30,1 & 9.630 & 35,9 & 1,19 \\
12 & parto normal; primeira gestação: sim & 145.427 & 13,9 & 5.073 & 18,9 & 1,36 \\
13 & sexo: M; despesa total: 0--185 & 218.959 & 20,9 & 6.842 & 25,5 & 1,22 \\
14 & parto normal; despesa total: 0--185 & 269.357 & 25,7 & 7.999 & 29,8 & 1,16 \\
15 & sexo: M; primeira gestação: sim & 135.847 & 13,0 & 4.562 & 17,0 & 1,31 \\

\secaotabela{População indígena}
16 & nascimento domiciliar & 1.593 & 4,3 & 108 & 12,5 & 2,94 \\
17 & primeira gestação: sim & 7.616 & 20,3 & 229 & 26,5 & 1,30 \\
18 & parto assistido por: outros & 3.511 & 9,4 & 134 & 15,5 & 1,66 \\
19 & domicílio particular improvisado & 2.485 & 6,6 & 107 & 12,4 & 1,87 \\
20 & cômodos no domicílio: 1--2; nascimento domiciliar & 603 & 1,6 & 60 & 7,0 & 4,32 \\

\secaotabela{População quilombola}
21 & famílias no domicílio: 1 ou menos; material do piso: cerâmica/pedra/carpete; destino do lixo: não coletado; banheiro: não & 4.743 & 21,9 & 205 & 29,0 & 1,32 \\
22 & famílias no domicílio: 1 ou menos; material do piso: cerâmica/pedra/carpete; destino do lixo: não coletado; abastecimento de água: poço/nascente & 5.950 & 27,5 & 243 & 34,4 & 1,25 \\
23 & famílias no domicílio: 1 ou menos; material do piso: cerâmica/pedra/carpete; destino do lixo: não coletado; material do piso: outros & 4.317 & 20,0 & 188 & 26,6 & 1,33 \\
24 & destino do lixo: não coletado; material do piso: cerâmica/pedra/carpete; despesa com energia: 0 & 4.151 & 19,2 & 175 & 24,8 & 1,29 \\
25 & material do piso: cerâmica/pedra/carpete; famílias no domicílio: 1 ou menos & 7.253 & 33,5 & 274 & 38,8 & 1,16 \\

\secaotabela{Mães com baixa escolaridade}
26 & sexo: M & 98.930 & 50,9 & 2.818 & 54,8 & 1,08 \\
27 & consultas: 1--3 & 24.963 & 12,8 & 853 & 16,6 & 1,29 \\
28 & destino do lixo: não coletado & 64.703 & 33,3 & 1.871 & 36,4 & 1,09 \\
29 & sexo: M; raça/cor da mãe: preta/parda; calçamento do domicílio: não existe & 39.468 & 20,3 & 1.189 & 23,1 & 1,14 \\
30 & sexo: M; despesa total: 0--185 & 29.614 & 15,2 & 917 & 17,8 & 1,17 \\

\secaotabela{Mães menores de 15 anos}
31 & consultas: 1--3 & 2.884 & 11,3 & 168 & 20,9 & 1,86 \\
32 & parto normal; sexo: M & 7.333 & 28,6 & 303 & 37,8 & 1,32 \\
33 & parto normal & 14.836 & 57,9 & 512 & 63,8 & 1,10 \\
34 & parto normal; sexo: M; consultas: 1--3 & 994 & 3,9 & 78 & 9,7 & 2,51 \\
35 & sexo: M; despesa com água/esgoto: 0 & 5.757 & 22,5 & 226 & 28,2 & 1,25 \\

\secaotabela{Beneficiárias do Programa Bolsa Família}
36 & sexo: M; raça/cor da mãe: preta/parda & 407.456 & 34,9 & 10.199 & 40,3 & 1,15 \\
37 & sexo: M; calçamento do domicílio: não existe & 268.135 & 23,0 & 6.843 & 27,0 & 1,18 \\
38 & raça/cor da mãe: preta/parda; material do piso: cerâmica/pedra/carpete & 531.500 & 45,5 & 12.504 & 49,4 & 1,08 \\
39 & destino do lixo: não coletado & 296.909 & 25,4 & 7.370 & 29,1 & 1,14 \\
40 & primeira gestação: sim & 159.938 & 13,7 & 4.337 & 17,1 & 1,25 \\
\bottomrule
\end{tabularx}
%\fontetabela
\end{table}

%Geral
%No recorte geral, o sexo masculino aparece em quatro dos cinco subgrupos (1, 3, 4 e 5), o parto normal em dois (2 e 3) e a primeira gestação em dois (2 e 5).

%O subgrupo 2 (parto normal e primeira gestação) apresenta o maior risco do desfecho da tabela (36\% acima da basal). O subgrupo 1, definido apenas pelo sexo masculino, tem o menor risco do desfecho (10\% acima da basal) e a maior participação nas ocorrências (56,2\%); a menor participação nas ocorrências é a do subgrupo 5 (17,2\%).

%O subgrupo 1 permite duas comparações com um item de diferença: a adição da ausência de calçamento eleva o risco de 10\% para 20\% acima da basal (subgrupo 4), e a adição da primeira gestação, para 29\% (subgrupo 5), com queda da participação nas ocorrências de 56,2\% para 23,4\% e 17,2\%, respectivamente.

%10k
%Nos municípios de até 10 mil habitantes, o parto normal aparece em três subgrupos (5, 8 e 10) e o sexo masculino em dois (6 e 9). O número de consultas aparece em um subgrupo, na faixa de 4 a 6 (subgrupo 7).

%O subgrupo 8 (parto normal e primeira gestação) apresenta o maior risco do desfecho da tabela (40\% acima da basal) e a menor participação nas ocorrências (18,0\%). O subgrupo 10, definido apenas por parto normal, tem o menor risco do desfecho (8\% acima da basal) e a maior participação nas ocorrências (52,9\%).

%O subgrupo 10 permite duas comparações com um item de diferença: a adição do sexo masculino eleva o risco de 8\% para 24\% acima da basal (subgrupo 6), e a adição da primeira gestação, para 40\% (subgrupo 8), com queda da participação nas ocorrências de 52,9\% para 30,4\% e 18,0\%, respectivamente.

%Rural 
%Na população rural, os cinco subgrupos são pares formados a partir de quatro itens: parto normal (11, 12 e 14), sexo masculino (11, 13 e 15), primeira gestação (12 e 15) e despesa total da família entre 0 e 185 (13 e 14).

%O subgrupo 12 (parto normal e primeira gestação) apresenta o maior risco do desfecho da tabela (36\% acima da basal); a menor participação nas ocorrências é a do subgrupo 15 (17,0\%). O subgrupo 14 tem o menor risco do desfecho (16\% acima da basal), e o subgrupo 11 tem a maior participação nas ocorrências (35,9\%).

%Os dois subgrupos que contêm a primeira gestação (12 e 15) são os de maior risco (36\% e 31\% acima da basal) e os de menor participação nas ocorrências (18,9\% e 17,0\%).

%indigena 
%Na população indígena, quatro dos cinco subgrupos são definidos por um único item. O nascimento domiciliar aparece em dois (15 e 20); os demais itens descrevem a assistência ao parto (subgrupo 18) e as condições de moradia (subgrupos 19 e 20).

%O subgrupo 20 (nascimento domiciliar em domicílio de 1 a 2 cômodos) apresenta o maior risco do desfecho da tabela (332\% acima da basal, isto é, 4,32 vezes a basal), com prevalência de 9,95\% e a menor participação nas ocorrências (7,0\%). O subgrupo 16, definido apenas pelo nascimento domiciliar, está 194\% acima da basal, com prevalência de 6,78\%. O subgrupo 17, definido apenas por primeira gestação, tem o menor risco do desfecho (30\% acima da basal) e a maior participação nas ocorrências (26,5\%).

%Os subgrupos 16 e 20 diferem por um único item: a adição de 1 a 2 cômodos no domicílio ao nascimento domiciliar eleva o risco de 194\% para 332\% acima da basal, com queda da participação nas ocorrências de 12,5\% para 7,0\%.

%Quilombola
%Na população quilombola, o material do piso (cerâmica, pedra ou carpete) aparece nos cinco subgrupos, o destino do lixo não coletado em quatro (21 a 24) e famílias no domicílio (1 ou menos) também em quatro (21, 22, 23 e 25). Os subgrupos 21 a 23 compartilham três itens e diferem no quarto: ausência de banheiro, abastecimento de água por poço ou nascente e material do piso.

%Os riscos variam pouco, entre 16\% e 33\% acima da basal. O subgrupo 23 apresenta o maior risco do desfecho (33\% acima da basal), seguido de perto pelo subgrupo 21 (32\%). O subgrupo 25, que reúne apenas material do piso e famílias no domicílio, tem o menor risco do desfecho (16\% acima da basal) e a maior participação nas ocorrências (38,8\%); a menor participação é a do subgrupo 24 (24,8\%).

%mães com baixa escolaridade
%Nas mães com baixa escolaridade, o sexo masculino aparece em três subgrupos (26, 29 e 30) e os demais itens ocorrem uma vez cada. O número de consultas aparece em um subgrupo, na faixa de 1 a 3 (subgrupo 27).

%O subgrupo 27, definido apenas pela faixa de 1 a 3 consultas, apresenta o maior risco do desfecho da tabela (29\% acima da basal) e a menor participação nas ocorrências (16,6\%). O subgrupo 26, definido apenas pelo sexo masculino, tem o menor risco do desfecho (8\% acima da basal) e a maior participação nas ocorrências (54,8\%).

%Os subgrupos 26 e 30 diferem por um único item: a adição da despesa total da família entre 0 e 185 eleva o risco de 8\% para 17\% acima da basal, com queda da participação nas ocorrências de 54,8\% para 17,8\%.

%Maes < 15
%Nas mães menores de 15 anos, o parto normal aparece em três subgrupos (32, 33 e 34), o sexo masculino em três (32, 34 e 35) e a faixa de 1 a 3 consultas em dois (31 e 34).

%O subgrupo 34 (parto normal, sexo masculino e 1 a 3 consultas) apresenta o maior risco do desfecho da tabela (151\% acima da basal, isto é, 2,51 vezes a basal), com prevalência de 7,85\% e a menor participação nas ocorrências (9,7\%). O subgrupo 33, definido apenas por parto normal, tem o menor risco do desfecho (10\% acima da basal) e a maior participação nas ocorrências (63,8\%).

%Os subgrupos 33, 32 e 34 formam uma sequência com um item de diferença a cada passo: o parto normal isolado está 10\% acima da basal; com o sexo masculino, 32\%; e com a faixa de 1 a 3 consultas, 151\%. A faixa de 1 a 3 consultas isolada (subgrupo 1) está 86\% acima da basal.

%BF
%Nos Beneficiárias do Bolsa Família, o sexo masculino e a raça/cor da mãe preta ou parda aparecem em dois subgrupos cada (36 e 37; 36 e 38), e os demais itens ocorrem uma vez cada.

%O subgrupo 40, definido apenas por primeira gestação, apresenta o maior risco do desfecho da tabela (25\% acima da basal) e a menor participação nas ocorrências (17,1\%). O subgrupo 38 tem o menor risco do desfecho (8\% acima da basal) e a maior participação nas ocorrências (49,4\%). Todos os subgrupos ficam até 25\% acima da basal.

%todos os recortes
%No Apgar baixo no quinto minuto, os itens mais recorrentes descrevem o recém-nascido e o parto. O sexo masculino aparece em seis dos oito recortes, ausente apenas nas populações indígena e quilombola; isolado, está entre 8\% e 10\% acima da basal (mães com baixa escolaridade e recorte geral). O parto normal aparece em quatro recortes (recorte geral, municípios de até 10 mil habitantes, rural e mães menores de 15 anos), e a conjunção de parto normal e primeira gestação se repete em três deles com risco semelhante: 36\% no recorte geral, 40\% nos municípios de até 10 mil habitantes e 36\% no rural. A primeira gestação isolada está 30\% acima da basal na população indígena e 25\% entre as beneficiárias do Bolsa Família.

Diferentemente do observado nos desfechos anteriores, os itens de infraestrutura que aparecem no Apgar baixo descrevem, em sua maioria, condições ausentes ou precárias: ausência de calçamento (recorte geral, mães com baixa escolaridade e beneficiárias do Bolsa Família), lixo não coletado (quilombola, mães com baixa escolaridade e beneficiárias do Bolsa Família), ausência de banheiro e abastecimento por poço ou nascente (quilombola), domicílio improvisado ou com 1 a 2 cômodos (indígena) e despesa total da família entre 0 e 185 (rural e mães com baixa escolaridade). A raça/cor da mãe preta ou parda aparece no recorte geral, entre mães com baixa escolaridade e entre beneficiárias do Bolsa Família.

O número de consultas de pré-natal aparece em três recortes: a faixa de 1 a 3 consultas entre mães com baixa escolaridade (29\% acima da basal) e entre mães menores de 15 anos (86\% isolada e 151\% combinada a parto normal e sexo masculino), e a faixa de 4 a 6 consultas nos municípios de até 10 mil habitantes (19\%).

Os extremos estão em dois recortes. A população indígena reúne os dois maiores riscos de todo o desfecho, ambos ligados ao nascimento domiciliar: 332\% acima da basal quando combinado a 1 a 2 cômodos (prevalência de 9,95\%) e 194\% isoladamente (prevalência de 6,78\%). Em seguida estão as mães menores de 15 anos, com 151\% acima da basal (prevalência de 7,85\%). Nos outros seis recortes, o subgrupo de maior risco fica entre 25\% (beneficiárias do Bolsa Família) e 40\% (municípios de até 10 mil habitantes) acima da basal.

\FloatBarrier
